\section{Cosmology and the Dark Sector (P10, P13, P18, P19, P29, P30)}
\label{sec:cosmology}

\subsection{The arrow of time and emergent expansion (P13)}

\textbf{Where we landed:} classical sequential growth gives a clean arrow of time, the
number of causal relations accumulates monotonically and irreversibly, the irreversible
cascade of the bootstrap. Expansion is demonstrated two ways. Sprinkled into a de Sitter
geometry the causal order expands (its intrinsic spatial slices grow with proper time).
And, more deeply, expansion EMERGES from a pure local growth rule: with net birth above
one (a second child with probability $q$) the spatial front grows on its own at rate
$1+q$, with a static control at $q=0$, and no metric imposed. The remaining frontier is
a fully emergent spatial geometry, not just an emergent growth rate.

\subsection{Inflation (P30)}

\textbf{Where we landed:} a time-varying birth rate gives the inflationary profile. A
high early rate produces a burst of rapid expansion ($4.2$ e-folds, the front doubling
each generation) followed by a graceful exit to slow expansion (rate about $1.06$), all
from a local rule with no inflaton field put in by hand.

\subsection{Dark energy: the everpresent cosmological constant (P10, P19, P29)}

The everpresent model: $\Lambda$ is conjugate to the spacetime $4$-volume $V$, realized
as the element count, which fluctuates as $\sqrt{V}$, so $\delta\Lambda \sim 1/\sqrt{V}$.
\textbf{Where we landed:} the sharp Benincasa-Dowker action has the fluctuation problem
(the implied $\Lambda$ grows with volume), and smearing tames it. In 2D the smeared
action makes the implied $\Lambda$ shrink with volume (exponent $-0.29$), the everpresent
direction. In 4D the (newly implemented) smeared kernel tames the exponent from $+0.16$
toward zero ($+0.06$), the same direction, though not yet the full shrinking, which
likely needs the dynamical conjugate-volume model. The observational payoff is in
Section~\ref{sec:predictions}.

\subsection{Dark matter (P18)}

\textbf{Where we landed:} a nonlocal (infrared-enhanced) gravitational kinetic term,
natural because the discrete action is nonlocal, flattens galaxy rotation curves with no
dark particle. The rotation curve falls under local (Newtonian) gravity (outer-to-inner
$v^2$ ratio $0.23$) but stays flat or rises under the nonlocal one (ratio $1.29$). This
is the mechanism and the flattening, demonstrated, the galactic scale being a separate
quantitative question. Two further routes (a gauge-neutral fermion that gravitates but
is electromagnetically dark, and relics from the early growth) remain.
